\subsection{Activation Checkpointing（梯度检查点）}

\begin{importantbox}{Activation Checkpointing 的核心思想}
反向传播需要前向过程中保存的激活值来计算梯度。对于深层网络，所有层的激活值同时驻留在 GPU 内存中，内存开销巨大。

\textbf{Activation checkpointing}（也叫 gradient checkpointing 或 rematerialization）的策略是：只保存部分层的激活值（称为 checkpoint），丢弃其余层的激活值。反向传播到某一层时，从最近的 checkpoint 重新做前向计算来恢复所需的激活值。

这是一个经典的\textbf{时间换空间}权衡。
\end{importantbox}

假设模型有 $L$ 层，每层激活值占 $a$ 内存：

\begin{table}[H]
\centering
\begin{tabular}{p{0.22\textwidth}p{0.22\textwidth}p{0.22\textwidth}p{0.18\textwidth}}
\toprule
\textbf{策略} & \textbf{激活值内存} & \textbf{额外计算} & \textbf{适用场景} \\
\midrule
不做 checkpointing & $L \cdot a$ & 0 & 内存充足 \\
每 $\sqrt{L}$ 层 checkpoint & $\sqrt{L} \cdot a$ & $\sim 33\%$ 前向 & 经典权衡 \\
每层 checkpoint & $a$ (仅 1 层) & $\sim 100\%$ 前向 & 极端内存受限 \\
\bottomrule
\end{tabular}
\caption{Activation checkpointing 的不同策略}
\end{table}

PyTorch 提供了 \texttt{torch.utils.checkpoint.checkpoint} 函数来实现这一功能：

\begin{lstlisting}[caption={使用 activation checkpointing 的示例}]
from torch.utils.checkpoint import checkpoint

class TransformerBlock(nn.Module):
    def forward(self, x):
        # ... transformer block logic ...
        return output

class Model(nn.Module):
    def __init__(self, num_layers):
        super().__init__()
        self.layers = nn.ModuleList(
            {[}TransformerBlock() for _ in range(num_layers){]}
        )

    def forward(self, x):
        for layer in self.layers:
            # Wrap each layer with checkpointing
            x = checkpoint(layer, x, use_reentrant=False)
        return x
\end{lstlisting}

\begin{warningbox}{Activation Checkpointing 的陷阱}
\begin{itemize}
    \item \textbf{计算开销}：每个被 checkpoint 的段在反向时会重新执行前向计算，总计算量增加约 33\%（如果每 $\sqrt{L}$ 层 checkpoint）。
    \item \textbf{不可 checkpoint 的操作}：含有副作用（如 dropout 的随机状态）的层需要特殊处理，否则重新计算时行为会不一致。
    \item \textbf{与 AMP 的交互}：在混合精度训练中使用 checkpointing 时，需要确保 autocast 上下文在 checkpoint 段内正确传播。
\end{itemize}
\end{warningbox}

\begin{knowledgebox}{Activation Checkpointing 的内存收益实例}
以 70B 模型（约 80 层 Transformer）为例，假设每层激活值约占 200 MB（取决于 batch size 和序列长度）：
\begin{itemize}
    \item 不做 checkpointing：$80 \times 200 = 16,000$ MB $= 16$ GB
    \item 每 $\sqrt{80} \approx 9$ 层 checkpoint：$9 \times 200 = 1,800$ MB $\approx 1.8$ GB
    \item 代价：约 33\% 额外前向计算
\end{itemize}
这意味着可以节省约 14 GB 激活值内存，代价是每步训练时间增加约 33\%。在 GPU 内存紧张时，这个权衡几乎总是值得的。
\end{knowledgebox}

